% =========================================================================
%  Ecuaciones Diferenciales — Tema 5: EDO lineales de 2.º orden con
%  coeficientes variables. Parte I: puntos ordinarios y series enteras.
%
%  Estilo: artículo matemático sobre fondo oscuro (texto claro en Computer
%  Modern, como en las láminas de análisis complejo de referencia), con
%  enunciados numerados por sección, cajas con pestaña de título para
%  definiciones y teoremas y figuras TikZ/pgfplots adaptadas al fondo negro.
% =========================================================================

% --- IDIOMA Y CODIFICACIÓN ------------------------------------------------
% Se compila con LuaLaTeX (memoria dinámica para las superficies 3D);
% con pdfLaTeX funciona igualmente reduciendo la resolución de esas figuras.
\usepackage{iftex}
\ifPDFTeX
  \usepackage[utf8]{inputenc}
  \usepackage[T1]{fontenc}
  \usepackage{lmodern}
\else
  \usepackage{fontspec}   % Latin Modern (Computer Modern OpenType) por defecto
\fi
\usepackage[spanish,es-tabla,es-nodecimaldot,es-noshorthands,es-lcroman]{babel}

% --- TIPOGRAFÍA (Latin Modern = Computer Modern vectorial) ----------------
\usepackage[final]{microtype}
\linespread{1.06}
\emergencystretch=2em

% --- MATEMÁTICAS -------------------------------------------------------------
\usepackage{amsmath, amssymb, mathtools}
\usepackage{bm}
\usepackage{mathrsfs}
\usepackage{cancel}

% --- PÁGINA -----------------------------------------------------------------
\usepackage[a4paper, top=2.9cm, bottom=2.7cm, left=2.7cm, right=2.7cm,
            headsep=0.9cm, headheight=22pt]{geometry}
\usepackage{graphicx}

% --- PALETA OSCURA ------------------------------------------------------------
\usepackage[table]{xcolor}
\definecolor{fondo}{HTML}{000000}      % fondo de página (negro puro)
\definecolor{texto}{HTML}{ECECEC}      % texto principal (blanco suave)
\definecolor{blanco}{HTML}{FFFFFF}
\definecolor{negro}{HTML}{000000}
\definecolor{gris}{HTML}{9CA0A8}       % texto secundario, ejes
\definecolor{grismedio}{HTML}{5B5F67}  % reglas, marcos tenues
\definecolor{grisoscuro}{HTML}{2E3036} % rejillas, divisiones
\definecolor{cajafondo}{HTML}{0E1014}  % interior de las cajas
\definecolor{cian}{HTML}{20B4CE}       % definiciones y acento principal
\definecolor{cianclaro}{HTML}{7FDBEA}  % enlaces
\definecolor{violeta}{HTML}{A48BFA}    % teoremas, proposiciones, lemas
\definecolor{ambar}{HTML}{F4B740}      % resultados clave, receta
\definecolor{magenta}{HTML}{E5579D}
\definecolor{verde}{HTML}{4FD68A}
\definecolor{naranja}{HTML}{FA8F45}
\definecolor{rojo}{HTML}{F25F5C}

% Página negra y texto claro. Se fija en el preámbulo para que el color
% «por defecto» (el que restauran flotantes, notas, cabeceras...) sea el claro.
\pagecolor{fondo}
\color{texto}

% --- GRÁFICOS -----------------------------------------------------------------
\usepackage{tikz}
\usetikzlibrary{arrows.meta, calc, positioning, decorations.pathmorphing,
                decorations.markings, decorations.pathreplacing, patterns,
                shapes.geometric, shapes.misc, fit, backgrounds, shadings,
                fadings, intersections}
\usepackage{pgfplots}
\pgfplotsset{compat=1.18}
\usepgfplotslibrary{colormaps, fillbetween, groupplots}

% --- LISTAS, TABLAS, CABECERAS ----------------------------------------------
\usepackage{booktabs}
\usepackage{array}
\usepackage{multirow}
\usepackage{enumitem}
\setlist[itemize]{leftmargin=1.6em, itemsep=2pt, topsep=4pt,
                  label={\color{cian}\scriptsize$\blacktriangleright$}}
\setlist[enumerate]{leftmargin=1.9em, itemsep=2pt, topsep=4pt}
\arrayrulecolor{grismedio}

\usepackage{fancyhdr}
\pagestyle{fancy}
\fancyhf{}
\fancyhead[LE]{\small\color{gris}\thepage\hfill\scshape Ecuaciones Diferenciales}
\fancyhead[RO]{\small\color{gris}\scshape Tema 5\,·\,Puntos ordinarios y series enteras\hfill\upshape\thepage}
\fancyhead[LO,RE]{}
\renewcommand{\headrulewidth}{0.3pt}
\renewcommand{\headrule}{\hbox to\headwidth{\color{grisoscuro}\leaders\hrule height \headrulewidth\hfill}}
\fancypagestyle{plain}{\fancyhf{}\renewcommand{\headrulewidth}{0pt}%
  \fancyfoot[C]{\small\color{gris}\thepage}}

% --- TÍTULOS DE SECCIÓN -----------------------------------------------------
\usepackage{titlesec}
\titleformat{\section}[block]
  {\normalfont\Large\bfseries\color{blanco}}
  {\color{cian}\thesection}{0.9em}{}
  [{\vspace{3pt}\color{grisoscuro}\titlerule[0.5pt]}]
\titlespacing*{\section}{0pt}{4.2ex plus 1ex minus .4ex}{2.2ex plus .4ex}
\titleformat{\subsection}[block]
  {\normalfont\large\bfseries\color{blanco}}
  {\color{cian}\thesubsection}{0.8em}{}
\titlespacing*{\subsection}{0pt}{3.2ex plus 1ex minus .4ex}{1.4ex plus .3ex}
\titleformat{\paragraph}[runin]
  {\normalfont\bfseries\color{blanco}}{}{0pt}{}[.]
\titlespacing*{\paragraph}{0pt}{1.6ex plus .5ex}{0.6em}
% Evita títulos huérfanos al pie de página
\usepackage{needspace}
\usepackage{placeins}
\usepackage{etoolbox}
\AtBeginDocument{%
  \pretocmd{\section}{\needspace{9\baselineskip}}{}{}%
  \pretocmd{\subsection}{\needspace{6\baselineskip}}{}{}}
\setcounter{secnumdepth}{2}
\setcounter{tocdepth}{2}

% --- ÍNDICE -------------------------------------------------------------------
\usepackage{tocloft}
\renewcommand{\cfttoctitlefont}{\normalfont\Large\bfseries\color{blanco}}
\renewcommand{\cftsecfont}{\bfseries\color{blanco}}
\renewcommand{\cftsecpagefont}{\bfseries\color{cian}}
\renewcommand{\cftsubsecfont}{\color{texto}}
\renewcommand{\cftsubsecpagefont}{\color{gris}}
\renewcommand{\cftsecleader}{\color{grismedio}\cftdotfill{\cftdotsep}}
\renewcommand{\cftsubsecleader}{\color{grisoscuro}\cftdotfill{\cftdotsep}}
\setlength{\cftbeforesecskip}{6pt}
\renewcommand{\contentsname}{Índice}

% --- ENLACES Y REFERENCIAS CRUZADAS ----------------------------------------
\usepackage[colorlinks=true, linkcolor=cianclaro, citecolor=ambar,
            urlcolor=cianclaro, plainpages=false, pdfusetitle=false,
            pdftitle={Tema 5. Puntos ordinarios y series enteras},
            pdfauthor={Daniel Martín Guerrero},
            pdfsubject={Ecuaciones Diferenciales},
            bookmarksnumbered=true]{hyperref}
\usepackage{amsthm}
\usepackage[spanish,nameinlink]{cleveref}   % tras amsthm y hyperref

% --- BIBLIOGRAFÍA -----------------------------------------------------------
\usepackage{csquotes}
\usepackage[backend=biber, style=numeric, sorting=none, giveninits=true,
            maxbibnames=6, minbibnames=3, doi=false, url=false,
            isbn=false]{biblatex}
\addbibresource{referencias.bib}
\renewcommand*{\bibfont}{\small}
\setlength{\bibitemsep}{4pt}

% --- PIES DE FIGURA -------------------------------------------------------
\usepackage[font={small,color=texto}, labelfont={sc,color=cian},
            labelsep=period, margin=0.4cm, justification=justified,
            skip=8pt]{caption}
\usepackage{subcaption}
\numberwithin{equation}{section}
\numberwithin{figure}{section}
\numberwithin{table}{section}

% =========================================================================
%               NOTACIÓN Y MACROS
% =========================================================================
\newcommand{\dif}{\mathrm{d}}
\newcommand{\eul}{\mathrm{e}}
\newcommand{\iu}{\mathrm{i}}
\newcommand{\R}{\mathbb{R}}
\newcommand{\C}{\mathbb{C}}
\newcommand{\N}{\mathbb{N}}
\newcommand{\Z}{\mathbb{Z}}
\newcommand{\abs}[1]{\lvert #1\rvert}
\newcommand{\Abs}[1]{\left\lvert #1\right\rvert}
\newcommand{\defeq}{\mathrel{\overset{\mathrm{def}}{=}}}
\newcommand{\Cinf}{\mathcal{C}^{\infty}}
\newcommand{\Com}{\mathcal{C}^{\omega}}           % funciones analíticas reales
\newcommand{\disco}[2]{D(#1,#2)}                  % disco abierto
\newcommand{\dd}[2]{\frac{\dif #1}{\dif #2}}
\newcommand{\ddd}[2]{\frac{\dif^{2} #1}{\dif #2^{2}}}
\newcommand{\serie}[3][n]{\sum_{#1=#2}^{\infty}#3}
\DeclareMathOperator{\sen}{sen}
\DeclareMathOperator{\senh}{senh}
\DeclareMathOperator{\artanh}{artanh}
\DeclareMathOperator{\dist}{dist}
\DeclareMathOperator{\Real}{Re}
\DeclareMathOperator{\Imag}{Im}
\DeclareMathOperator{\Arg}{arg}
\DeclareMathOperator{\ord}{ord}
\DeclareMathOperator{\Ai}{Ai}
\DeclareMathOperator{\Bi}{Bi}
\renewcommand{\sin}{\sen}
\newcommand{\Heav}{\Theta}
\newcommand{\nota}[1]{\textcolor{gris}{#1}}        % texto secundario
\newcommand{\resalta}[1]{\textcolor{blanco}{\bfseries #1}}
% Llaves horizontales en color tenue
\newcommand{\ubr}[2]{\underbrace{#1}_{\textcolor{gris}{#2}}}

% =========================================================================
%        ENTORNOS FORMALES
% =========================================================================
\usepackage[most]{tcolorbox}

% Contador común a todos los enunciados (numeración «sección.n»)
\newcounter{enunciado}[section]
\renewcommand{\theenunciado}{\thesection.\arabic{enunciado}}

% --- Cajas con pestaña (definiciones, teoremas, resultados) -------------------
% #1 color del marco, #2 rótulo completo
\tcbset{
  caja/.style={enhanced, colback=cajafondo, colupper=texto,
    colframe=#1, boxrule=0.7pt, arc=2.5pt, outer arc=2.5pt,
    left=10pt, right=10pt, top=11pt, bottom=8pt,
    before skip=16pt plus 2pt, after skip=14pt plus 2pt,
    attach boxed title to top left={yshift*=-\tcboxedtitleheight/2, xshift=10pt},
    boxed title style={colback=#1, colframe=#1, arc=2pt, outer arc=2pt,
                       boxrule=0pt, left=5pt, right=5pt, top=1.5pt, bottom=1.5pt},
    fonttitle=\sffamily\bfseries\small, coltitle=negro},
}
% Constructor: \cajaenunciado{nombre}{Rótulo}{color}{estilo cuerpo}
\makeatletter
\newcommand{\cajaenunciado}[4]{%
  \newcounter{#1}%
  \expandafter\def\csname the#1\endcsname{\theenunciado}%
  \expandafter\def\csname theH#1\endcsname{#1.\arabic{#1}}%
  \newtcolorbox{tcb@#1}[1]{caja=#3,
      title={#2~\theenunciado\ifx\relax##1\relax\else\ \textnormal{\sffamily(##1)}\fi}}%
  \newenvironment{#1}[1][]%
    {\par\stepcounter{enunciado}\refstepcounter{#1}%
     \begin{tcb@#1}{##1}#4\ignorespaces}%
    {\end{tcb@#1}}%
}
\makeatother
\cajaenunciado{definicion}{Definición}{cian}{\normalfont}
\cajaenunciado{teorema}{Teorema}{violeta}{\itshape}
\cajaenunciado{proposicion}{Proposición}{violeta}{\itshape}
\cajaenunciado{lema}{Lema}{violeta}{\itshape}
\cajaenunciado{corolario}{Corolario}{violeta}{\itshape}
\cajaenunciado{resultado}{Resultado}{ambar}{\normalfont}

% Caja de método (receta), sin numerar
\newtcolorbox{metodo}[1][Método]{caja=ambar, title={#1}}

% --- Enunciados sin caja (ejemplos, observaciones, notación) -----------------
\newtheoremstyle{ejemplo}{12pt}{10pt}{\normalfont}{}%
  {\bfseries\color{cian}}{.}{0.6em}%
  {\thmname{#1}\thmnumber{ #2}\thmnote{\normalfont\color{gris}\ (#3)}}
\newtheoremstyle{observacion}{10pt}{10pt}{\normalfont}{}%
  {\itshape\bfseries\color{violeta}}{.}{0.6em}%
  {\thmname{#1}\thmnumber{ #2}\thmnote{\normalfont\itshape\color{gris}\ (#3)}}
\theoremstyle{ejemplo}
\newtheorem{ejemplo}[enunciado]{Ejemplo}
\newtheorem{notacion}[enunciado]{Notación}
\theoremstyle{observacion}
\newtheorem{observacion}[enunciado]{Observación}

% Demostraciones
\renewcommand{\qedsymbol}{\textcolor{gris}{$\blacksquare$}}
\makeatletter
\renewenvironment{proof}[1][\proofname]{\par
  \pushQED{\qed}\normalfont\topsep6\p@\@plus6\p@\relax
  \trivlist\item[\hskip\labelsep\itshape\bfseries\color{gris}#1\@addpunct{.}]\ignorespaces
}{\popQED\endtrivlist\@endpefalse}
\makeatother
\newenvironment{esbozo}{\begin{proof}[Esbozo de la demostración]}{\end{proof}}

% Notas al margen del discurso: «idea», recuadro tenue con barra lateral
\newtcolorbox{idea}[1][Idea]{enhanced, breakable, colback=cajafondo, colupper=texto,
  frame hidden, borderline west={2pt}{0pt}{cian!80}, arc=0pt,
  left=10pt, right=8pt, top=6pt, bottom=6pt, before skip=12pt, after skip=12pt,
  fontupper=\normalfont,
  before upper={{\sffamily\bfseries\small\color{cian}#1.}\ }}

\crefname{definicion}{definición}{definiciones}
\Crefname{definicion}{Definición}{Definiciones}
\crefname{teorema}{teorema}{teoremas}
\Crefname{teorema}{Teorema}{Teoremas}
\crefname{proposicion}{proposición}{proposiciones}
\Crefname{proposicion}{Proposición}{Proposiciones}
\crefname{lema}{lema}{lemas}
\Crefname{lema}{Lema}{Lemas}
\crefname{corolario}{corolario}{corolarios}
\Crefname{corolario}{Corolario}{Corolarios}
\crefname{resultado}{resultado}{resultados}
\Crefname{resultado}{Resultado}{Resultados}
\crefname{ejemplo}{ejemplo}{ejemplos}
\Crefname{ejemplo}{Ejemplo}{Ejemplos}
\crefname{observacion}{observación}{observaciones}
\Crefname{observacion}{Observación}{Observaciones}
\crefname{notacion}{notación}{notaciones}
\Crefname{notacion}{Notación}{Notaciones}
\crefname{section}{sección}{secciones}
\Crefname{section}{Sección}{Secciones}
\crefname{subsection}{apartado}{apartados}
\Crefname{subsection}{Apartado}{Apartados}
\crefname{figure}{figura}{figuras}
\Crefname{figure}{Figura}{Figuras}
\crefname{table}{tabla}{tablas}
\Crefname{table}{Tabla}{Tablas}
\crefname{equation}{ecuación}{ecuaciones}
\Crefname{equation}{Ecuación}{Ecuaciones}
\crefformat{equation}{(#2#1#3)}
\Crefformat{equation}{Ecuación~(#2#1#3)}
\crefrangeformat{equation}{(#3#1#4)--(#5#2#6)}
\crefmultiformat{equation}{(#2#1#3)}{ y~(#2#1#3)}{, (#2#1#3)}{ y~(#2#1#3)}

% Modo rápido: si existe el fichero vacío «rapido.flag», las superficies y
% retratos de fase se generan con baja resolución (compilación en segundos).
\newif\ifrapido
\IfFileExists{rapido.flag}{\rapidotrue}{\rapidofalse}
\ifrapido
  \def\resPx{24}\def\resPy{34}\def\resF{41}
\else
  \def\resPx{85}\def\resPy{131}\def\resF{141}
\fi

% =========================================================================
%        ESTILOS GRÁFICOS COMUNES
% =========================================================================
% =========================================================================
%  Estilos gráficos comunes (TikZ / pgfplots) para fondo negro
% =========================================================================

% --- Mapa de colores «arcoíris» para coloreado por fase (dominio complejo) ---
% Recorre el círculo cromático completo: el color codifica arg f(z) en (-pi,pi].
\pgfplotsset{colormap={fase}{
  rgb255=(255,40,60) rgb255=(255,150,20) rgb255=(250,230,40)
  rgb255=(60,220,90) rgb255=(30,200,230) rgb255=(60,90,255)
  rgb255=(190,60,240) rgb255=(255,40,60)}}

% --- Estilos TikZ generales ------------------------------------------------
\tikzset{
  >=Stealth,
  every picture/.append style={line cap=round, line join=round},
  eje/.style={-{Stealth[length=5pt,width=4pt]}, line width=0.6pt, gris},
  ejeclaro/.style={-{Stealth[length=5pt,width=4pt]}, line width=0.7pt, texto},
  fino/.style={line width=0.4pt, grismedio},
  guia/.style={densely dashed, line width=0.45pt, grismedio},
  punteado/.style={densely dotted, line width=0.6pt, gris},
  etiq/.style={font=\small, text=texto},
  etiqgris/.style={font=\footnotesize, text=gris},
  polo/.style={cross out, draw=rojo, line width=1.1pt, minimum size=6pt, inner sep=0pt},
  punto/.style={circle, fill=blanco, inner sep=0pt, minimum size=4pt},
  puntocian/.style={circle, fill=cian, inner sep=0pt, minimum size=4.5pt},
  orienta/.style={decoration={markings, mark=at position #1 with
     {\arrow{Stealth[length=6pt,width=5pt]}}}, postaction={decorate}},
  orienta/.default=0.55,
  nodo/.style={draw=grismedio, fill=cajafondo, rounded corners=2pt,
     line width=0.5pt, inner xsep=6pt, inner ysep=4pt, font=\small, text=texto},
  nodocian/.style={nodo, draw=cian},
  nodoambar/.style={nodo, draw=ambar},
  nodovioleta/.style={nodo, draw=violeta},
  flecha/.style={-{Stealth[length=5pt,width=4.5pt]}, line width=0.6pt, gris},
}

% --- Estilo base de ejes pgfplots en oscuro ----------------------------------
\pgfplotsset{
  oscuro/.style={
    axis background/.style={fill=fondo},
    tick label style={font=\scriptsize, text=gris},
    label style={font=\small, text=texto},
    title style={font=\small, text=texto},
    axis line style={grismedio, line width=0.5pt},
    every tick/.style={grismedio},
    major tick length=3pt,
    grid=major, grid style={grisoscuro, line width=0.35pt},
    legend style={font=\scriptsize, draw=grisoscuro, fill=cajafondo,
                  text=texto, rounded corners=1.5pt, inner sep=3pt,
                  row sep=1pt},
    legend cell align=left,
    /pgf/number format/use comma,
    /pgf/number format/1000 sep={},
    unbounded coords=jump,
  },
  % Ejes cruzados por el origen (estilo «pizarra»)
  cruz/.style={oscuro, axis lines=middle, grid=none,
    axis line style={gris, line width=0.6pt,
      -{Stealth[length=5pt,width=4pt]}},
    xlabel style={anchor=west, font=\small},
    ylabel style={anchor=south, font=\small},
  },
  superficie/.style={
    tick label style={font=\scriptsize, text=gris},
    label style={font=\small, text=texto},
    axis line style={grismedio}, every tick/.style={grismedio},
    grid=none, colormap name=fase, shader=interp,
    /pgf/number format/use comma,
  },
  % Curvas
  c1/.style={line width=1.15pt, cian},
  c2/.style={line width=1.15pt, magenta},
  c3/.style={line width=1.15pt, ambar},
  c4/.style={line width=1.15pt, verde},
  c5/.style={line width=1.15pt, violeta},
  c6/.style={line width=1.15pt, naranja},
}


